
\documentclass{article}
\usepackage{colortbl}
\usepackage{makecell}
\usepackage{multirow}
\usepackage{supertabular}

\begin{document}

\newcounter{utterance}

\centering \large Interaction Transcript for game `cladder', experiment `full\_v1.5\_default', episode 1442 with Qwen3.6{-}27B.
\vspace{24pt}

{ \footnotesize  \setcounter{utterance}{1}
\setlength{\tabcolsep}{0pt}
\begin{supertabular}{c@{$\;$}|p{.15\linewidth}@{}p{.15\linewidth}p{.15\linewidth}p{.15\linewidth}p{.15\linewidth}p{.15\linewidth}}
    \# & \multicolumn{2}{c}{Player} && \multicolumn{2}{c}{Game Master} \\
    \hline

    \theutterance \stepcounter{utterance}  
    & & & \multicolumn{4}{p{0.6\linewidth}}{
        \cellcolor[rgb]{0.9,0.9,0.9}{
            \makecell[{{p{\linewidth}}}]{
                \texttt{\tiny{[P1$\langle$GM]}}
                \texttt{You are an expert in causal inference. The following question is not a typical commonsense query, but rather a meticulously designed question created by a professor specializing in causal inference, intended to assess the students' mastery of the course content.} \\
\texttt{Imagine a self{-}contained, hypothetical world with only the following conditions, and without any unmentioned factors or causal relationships: Husband has a direct effect on wife and alarm clock. Wife has a direct effect on alarm clock. The overall probability of alarm set by husband is 17\%. The probability of alarm not set by husband and ringing alarm is 69\%. The probability of alarm set by husband and ringing alarm is 3\%. Is the chance of ringing alarm smaller when observing alarm set by husband?} \\
\texttt{Start your answer with 'Yes' or 'No', followed by additional reasoning or evidence to support your explanation.} \\
            }
        }
    }
     \\ \\

    \theutterance \stepcounter{utterance}  
    & \multicolumn{4}{p{0.6\linewidth}}{
        \cellcolor[rgb]{0.561,0.737,0.561}{
            \makecell[{{p{\linewidth}}}]{
                \texttt{\tiny{[P1$\rangle$GM]}}
                \texttt{Yes} \\
\\ 
\texttt{**Reasoning:**} \\
\\ 
\texttt{To determine if the chance of a ringing alarm is smaller when observing that the alarm was set by the husband, we need to compare the conditional probability of the alarm ringing given that the husband set it, \$P(Ring | Husband=Set)\$, with the marginal probability of the alarm ringing, \$P(Ring)\$. Alternatively, we can compare \$P(Ring | Husband=Set)\$ with \$P(Ring | Husband=Not Set)\$.} \\
\\ 
\texttt{Let \$H\$ be the event that the husband sets the alarm.} \\
\texttt{Let \$R\$ be the event that the alarm rings.} \\
\\ 
\texttt{From the problem description, we have the following probabilities:} \\
\texttt{1.  \$P(H) = 0.17\$ (The overall probability of alarm set by husband).} \\
\texttt{2.  \$P(\textbackslash{}neg H \textbackslash{}cap R) = 0.69\$ (The probability of alarm not set by husband and ringing alarm).} \\
\texttt{3.  \$P(H \textbackslash{}cap R) = 0.03\$ (The probability of alarm set by husband and ringing alarm).} \\
\\ 
\texttt{First, let's calculate the marginal probability of the alarm ringing, \$P(R)\$:} \\
\texttt{\$\$P(R) = P(H \textbackslash{}cap R) + P(\textbackslash{}neg H \textbackslash{}cap R)\$\$} \\
\texttt{\$\$P(R) = 0.03 + 0.69 = 0.72\$\$} \\
\\ 
\texttt{Next, let's calculate the conditional probability} \\
            }
        }
    }
    & & \\ \\

    \theutterance \stepcounter{utterance}  
    & & & \multicolumn{2}{p{0.3\linewidth}}{
        \cellcolor[rgb]{0.9,0.9,0.9}{
            \makecell[{{p{\linewidth}}}]{
                \texttt{\tiny{[GM$|$GM]}}
                \texttt{Yes} \\
            }
        }
    }
    & & \\ \\

    \theutterance \stepcounter{utterance}  
    & & & \multicolumn{2}{p{0.3\linewidth}}{
        \cellcolor[rgb]{0.9,0.9,0.9}{
            \makecell[{{p{\linewidth}}}]{
                \texttt{\tiny{[GM$|$GM]}}
                \texttt{yes} \\
            }
        }
    }
    & & \\ \\

    \theutterance \stepcounter{utterance}  
    & & & \multicolumn{2}{p{0.3\linewidth}}{
        \cellcolor[rgb]{0.9,0.9,0.9}{
            \makecell[{{p{\linewidth}}}]{
                \texttt{\tiny{[GM$|$GM]}}
                \texttt{game\_result = WIN} \\
            }
        }
    }
    & & \\ \\

\end{supertabular}
}

\end{document}
